\cotefile{157-5}
\scanpage{453}\mspage{365}
% archive sheet number: 452
est \unsure{aussi|encore} bien dans $\mathrm{Fl}\,K'$.

Ainsi, $K'$ est une sous-catégorie de $J$, mais est-elle filtrante ? La condition $\mathrm{PS}\,1'$ est évidente, puisque $e$ est un plus petit élément de $J$, reste la condition $\mathrm{PS}\,2$ de la double flèche. On voit aisément que cette condition est satisfaite dans $K'$, sss elle est satisfaite pour les double-flèches
\diagram{\begin{tikzcd}[ampersand replacement=\&,column sep=large]
e \arrow[r,shift left=.48ex,"p^{\unsure{r|s}}"] \arrow[r,shift right=.48ex,"p^s"'] \& e \&[-1em] \text{avec}\quad r,s\geq 0\unsure{,|.}
\end{tikzcd}}

(\deleted{\unsure{ou encore}}En effet, \deleted{\unsure{des flèches|deux flèches}} On se ramène aussitôt au cas de \unsure{deux|ttes} flèches de but $e$, et on note que \unsure{tte} flèche de $K'$ de but $e$ s'écrit comme un composé
\diagram{\begin{tikzcd}[ampersand replacement=\&,column sep=small]
k\arrow[r,"u"] \& e\arrow[r,"p^r"] \& e\arrow[r,"v"] \& \unsure{e}
\end{tikzcd}}, avec $r\geq0$ et $u,v$
des flèches dans $K'$ \unsure{\ldots|.} En composant
avec \diagram{\begin{tikzcd}[ampersand replacement=\&,column sep=small]e\arrow[r,"p"] \& e\end{tikzcd}}, on trouve alors, comme $pv=p$
\unsure{:} \diagram{\begin{tikzcd}[ampersand replacement=\&,column sep=small]
k\arrow[r,"u"] \& e\arrow[r,"p^{r+1}"] \& e,\quad\text{soit}\quad \& k\arrow[r,"{\unsure{fu|pu}}"] \& e\arrow[r,"p^r"] \& e
\end{tikzcd}})

Mais on voit aussitôt que la possibilité d'égaliser $(p^r,p^s)$ dans $K'$ équivaut à l'existence de $t\unsure{>|\geq}0$ tel que $p^{r+t}=p^{s+t}$. Donc $K'$ est filtrante sss \deleted{\unsure{$\varphi$}} $r\longmapsto p^r$ est une \deleted{\unsure{suite}} stationnaire (faire $r\unsure{=|\geq}0$, $s=1$, d'où $p^t=p^{t+1}$ pour \unsure{grand|$t$ grand}). Ce n'est
\scanpage{454}\mspage{366}
% archive sheet number: 453
pas forcément le cas.

\deleted{Mais $\deleted{K}\unsure{J|K}$ \unsure{étant|est} \unsure{grande|gde} devant $\pi$,}

Posons donc \inserted{$K=K_0$,} $K'=\deleted{\unsure{K}}K_1$, et construisons par induction une suite croissante de sous-catégories $K_i$ ($i\in\unsure{\mathbb{N}|N}$), de $\mathrm{Ob}\,K_i=\mathrm{Ob}\,J$, et $K_i$ \inserted{(pour $i\geq1$)} étant déduite de $K_{i-1}$ en lui rajoutant \deleted{\unsure{des}} une flèche $p_i:e\longrightarrow e$, comme ci-dessus. Ainsi tte double flèche dans un $K_i$ peut s'égaliser dans $K_{i+1}$. Prenons $K_{\infty}^{\deleted{\unsure{\prime|q|,}}}=\unsure{\varinjlim|\underline{\lim}}\,K_i$, $K_{\infty}$ est une sous-catégorie \uline{filtrante} de $\deleted{K}J$, \unsure{contenant} $K$ et de \unsure{cardinal|card} $\leq\pi$, \deleted{\unsure{définissant une co-}} qed. $\big]$

On peut préciser cette construction ainsi\unsure{ ``|.}

\uline{Corollaire} 4.13.39. Sous les conditions de 4.13.38, supposons que $K$ contienne un objet $e$ \unsure{maximal|minimal} dans $\deleted{K}J$. Soit $\Sigma$ l'ens des endomorphismes $p\unsure{/}$de $e$ dans $J$, tels que
pour tte double-flèche \diagram{\begin{tikzcd}[ampersand replacement=\&,column sep=large]
k\arrow[r,shift left=.48ex,"u"]\arrow[r,shift right=.48ex,"v"'] \& e
\end{tikzcd}} \uline{dans} $K$, on
ait $pu=pv$. \deleted{Alors \unsure{a)} il existe \illegible{1},} Alors

\unsure{a)|c)} $\Sigma$ est stable par composition et il existe une partie non vide $\Sigma_{\infty}$ de $\Sigma$, telle que \deleted{pour \unsure{$p$} \illegible{1} \unsure{catégo}} pour \deleted{\unsure{tous}} $p,p'\in\Sigma_0$, existe
\scanpage{455}\mspage{367}
% archive sheet number: 454
$q\in\Sigma_0$ tel que $qp=qp'$ (i.e telle que la sous-catégorie \inserted{$J_0$} de $J$ dont l'ens des objets est $\{e\}$, et l'ens des flèches $\{\mathrm{id}_e\}\cup\Sigma_{\infty}$, soit filtrante).

Si $\Sigma_0$ est choisie ainsi, soit $F$ l'ens \deleted{des} \inserted{réunion de $\mathrm{Fl}\,\deleted{\unsure{J}}K$ et de l'ens des} flèches \deleted{de la \unsure{fo-}} composées
\diagram{\begin{tikzcd}[ampersand replacement=\&]
\unsure{(\alpha)|(a)} \& k\arrow[r,"\unsure{u}"] \& e\arrow[r,"p"] \& \deleted{\unsure{k}}e\arrow[r,"v"] \& k',
\end{tikzcd}}
avec $u,v\in\mathrm{Fl}\,J$, $p\in\Sigma_{\infty}$. Alors $F$ est l'ens. des flèches d'une \inserted{sous-}catégorie $\deleted{\unsure{J'}}K'$ de $J$, ayant \unsure{mes|m\^es|mêmes} objets que $J$, \deleted{et $K'$ est \unsure{fi-}} et admettant $e$ pour plus grand \unsure{élémt}. De plus, $K'$ contient $K$ et est filtrante.

\uline{\unsure{Dém}} Que $\Sigma$ soit stable par composition est clair. Pour construire $\Sigma_{\infty}$ ayant les propriétés énoncées, on note que, $\deleted{K}\unsure{J|K}$ \deleted{\unsure{étant}} \unsure{gde|grde} devant $\pi$, \deleted{et \unsure{f-} il existe} $\Sigma\deleted{\unsure{_0}}$ est non vide.$^*$
% slanted
\margin{cf argument plus bas pour \unsure{constr.|const.} de $p_i$} On construit \inserted{par \unsure{ordre|ind.}} une suite d'élémts
\[
p_0,p_1,\ldots\ p_n,\ldots
\]
de $\Sigma$, en prenant pour $p_0$ un élément quelconque de $\Sigma$, et pour $\unsure{i|j}\geq1$, prenant $p_i\in\Sigma$ tel que \inserted{pour} \unsure{tte} \deleted{\unsure{double}} double-flèche \deleted{\unsure{$(\alpha,\beta)$}} $(p,p')$ du \inserted{\unsure{ss-|s-}}monoïde \inserted{$M_{i-1}$} de $\mathrm{End}_J(e)$ engendré par $(p_j)_{0\leq j\leq i-1}$, \deleted{on} ait $p_i p=p_i p'$. Un tel $p_i$ existe, \deleted{puisque \unsure{par} \inserted{on \deleted{\unsure{remarque|puisque}}} \unsure{l'hyp.} \unsure{de} $J$ est \unsure{grand|grd} devant $\pi$ \unsure{donc} devant $\aleph_0$, et que $M_{i-1}$ est de cardinal $\leq\aleph_0$, il} en prenant l'ens des doubles-flèches de la forme précédente $(p,p')$, avec $p,p'\in M_{i-1}$, et de celles de la forme $(u,v):k\rightrightarrows e$, avec $u,v\in\mathrm{Fl}\,J$, l'ens des
\scanpage{456}\mspage{368}
% archive sheet number: 455
toutes ces doubles-flèches est de cardinal $\leq\pi$ ($M_{i-1}$ étant dénombrable, donc de cardinal $\leq\pi$), donc il existe une flèche $e\longrightarrow j$ dans $J$ qui les égalise toutes ($J$ étant \unsure{gde} devant $\pi$). Quitte à composer avec $j\longrightarrow e$ ($e$ étant objet maximal de $J$), on trouve le $p_i:e\longrightarrow e$ demandé.

Soit \unsure{alors|donc} $\deleted{\unsure{M_{\infty}|\Sigma_{\infty}}}\inserted{\unsure{\Sigma_0|\Sigma_{\infty}}}$ \unsure{le} \deleted{monoïde} \inserted{sous-ens. multiplicatif de $\Sigma_0$} engendré par les $p_i$ ($i\in\unsure{\mathbb{N}|N}$), \deleted{et $\Sigma_0$ \unsure{b-}} limite \unsure{inductive} des ens \unsure{similaires} $\Sigma_i$ \unsure{engendré|engendrés} par les $p_j$ pour $j\in[0,i]$. Si $p,p'$ sont dans $\Sigma_\infty$, ils sont dans un \inserted{$=$} $\Sigma_i$, donc on aura $p_{\unsure{1+i|j+i|i+1}}p=p_{\unsure{i+1|i-1}}p'$, ce qui prouve que $\Sigma_\infty$ a la propriété filtrante exigée.

Prouvons maintenant que $F$ est l'ens des flèches d'une sous-catégorie $\deleted{\unsure{J}}K'$ de $\deleted{K}\unsure{J|K}$, ayant $\mathrm{Ob}\,K$ comme ens d'objets. Il reste à établir la stabilité par composition $\unsure{vu|uv}$. Mais si $u$ et $v$, \deleted{\unsure{ou} $u$ \unsure{ou} $v$} sont dans $\mathrm{Fl}\,\deleted{\unsure{J}}K$, \deleted{cette \unsure{stabi-}} ou si l'une des deux est dans $\mathrm{Fl}\,K$, l'autre \unsure{non} (donc de la forme $\unsure{(\alpha)|(a)}$), cette stabilité est immédiate. Reste à voir le cas où \deleted{$u$ et $v$}\inserted{ni $u$ ni $v$} sont \deleted{tous deux} \inserted{dans} $\mathrm{Fl}\,K$, donc sont tous deux de la forme $(\alpha)$ :
\diagram{\begin{tikzcd}[ampersand replacement=\&,column sep=small]
k\arrow[r,"u"] \& e\arrow[r,"p"] \& e\arrow[r,"v"] \& k'\arrow[r,"u'"] \& e\arrow[r,"p'"] \& e\arrow[r,"\unsure{v'|u'}"] \& k''
\end{tikzcd}}
avec $u,v,u',v'\in\mathrm{Fl}\,K$, $p,p'\in\Sigma_{\infty}$. Comme $(u'v,\mathrm{id}_e)$ est une double-flèche de
\scanpage{457}\mspage{369}
% archive sheet number: 456
$K$ de but $e$, on a
\[
p'u'v=p'\mathrm{id}_e=p'
\]
donc le composé en question s'écrit aussi
\diagram{\begin{tikzcd}[ampersand replacement=\&]
k\arrow[r,"u"] \& e\arrow[r,"p'p"] \& e\arrow[r,"v'"] \& k'',
\end{tikzcd}}
où $p'p\in\Sigma_\infty$ puisque $\Sigma_\infty$ est stable par composition, donc il est dans \deleted{\unsure{$K'$}} $F$, qed.

Il est clair que $K'$ contient $K$, et ayant $\unsure{=\mathrm{s}|\Rightarrow}$ objets que $K$ qui admet $e$ comme plus grd élémt, que $e$ est un plus grd \unsure{élémt|élément} de $K'$. Reste à prouver que $K'$ est filtrante, Comme $K'$ a un plus grd élémt, il reste à voir \deleted{\illegible{1}} que toute double-flèche \inserted{$k\rightrightarrows l$} de $K'$ peut être égalisée dans $K'$. Quitte à composer avec $l\longrightarrow e$, \unsure{OPS.|OPL.} $l=e$. Or on \unsure{sait|voit} que toute flèche $k\longrightarrow e$ dans $K'$ est de la forme
\diagram{\begin{tikzcd}[ampersand replacement=\&]
k\arrow[r,"u"] \& e\arrow[r,"p"] \& e\arrow[r,"v"] \& \deleted{\unsure{k}}e,
\end{tikzcd}}
\deleted{et quitte à composer} avec $u,v\in\mathrm{Fl}\,K$ et $p\in\Sigma_\infty\cup\{\mathrm{id}_e\}$. Quitte à composer avec $p_0\in\Sigma_\infty$ \deleted{\unsure{choisi|fixé}}\unsure{fixé|choisi}, cette \deleted{\unsure{double}} flèche devient (puisque $p_0v=p_0\mathrm{id}_e=p_0$)
\diagram{\begin{tikzcd}[ampersand replacement=\&]
k\arrow[r,"u"] \& e\arrow[r,"p_0p"] \& e\unsure{..|:}
\end{tikzcd}}
\deleted{ou encore} i.e. $p_0pu$ \deleted{\unsure{pour notre}} Pour une flèche $v'p'u':k\longrightarrow e$, on \unsure{trouvera|trouve} $p_0p'u'$. Soit alors $q\in\Sigma_0$ tel que $q(p_0p)=q(p_0p')$, \inserted{soit $r$} ($q$\unsure{.} existe puisque $p_0p,p_0p'\in\Sigma_\infty$). Donc
\[
\big(q(p_0pu),q(\unsure{p_0^{.}|p_0'|p_0}p'u')\unsure{.}\big)=(r\unsure{u},ru'),
\]
et on a $ru=ru'$ puisque $r\in\Sigma_\infty$. Au total, $qp_0\deleted{\unsure{q}}$ égalise les deux flèches $vpu$ et $v'p'u'$, qed.
\scanpage{458}\mspage{370}
% archive sheet number: 457
\uline{4.14} : \uline{Critères d'accessibilité pour grosses catégories}.

\uline{Théorème} 4.14.1. Soient $C$ une catégorie équivalente à une petite catégorie, $\Phi,\inserted{\Psi}$ \deleted{deux} petits ens de \deleted{\unsure{foncteurs}} catégories, \deleted{\unsure{$\varphi:I\longrightarrow C$} Pour tte} équivalentes à des petites catégories. Pour tte catégorie $I$, on désigne par $I_\bullet$ (resp $ {}_\bullet I$) le cône droit (resp gauche) sur $I$, déduit de $I$ en \unsure{ajoutant} un objet final \unsure{$e_I$} (resp initial \inserted{$\deleted{\unsure{e}}e=e_I$}). Soit $\Lambda_1$ \inserted{(resp. $\Lambda_2$)} un ens. de couples $(\deleted{\unsure{\varphi}}I,\varphi)$ (resp. $(I,\psi)$), où $I\in\Phi$ (resp. $I\in\Psi$), et où
\begingroup\csname tagsleft@true\endcsname\[
\varphi:I_{\bullet}\longrightarrow C\qquad(\text{resp. }\psi:{}_\bullet I\longrightarrow C)
\tag*{(4.14.1.1)}
\]\endgroup
est un foncteur. Pour tte catégorie $\unsure{\mathcal N|N|\mathscr N}$ \unsure{stable}$^*$ par $\unsure{\varinjlim|\underline{\lim}}$ de type $\Phi$ et par $\varprojlim$ de type $\Psi$, on désigne par
% slanted
\margin{Cette propriété de stabilité n'est pas \unsure{néc.} pour la déf. \deleted{\unsure{naïve}} de $\underline{\mathrm{Hom}}_{\Lambda_1}^{\Lambda_2}(C,N)$, et il est \unsure{à peu près|toujours} sûr que \unsure{le th.} est \unsure{valable|applicable} sans faire cette hyp sur $N$. Cf 4.14.16, où \unsure{l'on|on ?} voit \unsure{que|ce} \unsure{pouvant|pour} \unsure{s'étendre|l'établir} \unsure{dans|en} \unsure{toute sa généralité|toute généralité} pour $\Psi$.}
\begingroup\csname tagsleft@true\endcsname\[
\underline{\mathrm{Hom}}_{\Lambda_1}^{\Lambda_2}(C,\mathcal N)
\subset\underline{\mathrm{Hom}}(C,\mathcal N)
\tag*{(4.14.1.2)}
\]\endgroup
la sous-catégorie \unsure{str} pleine de $\underline{\mathrm{Hom}}(C,\mathcal N)$ formée des foncteurs $f:C\longrightarrow\mathcal N$ \deleted{tel} ayant la propriété suivante :

a) $\forall\,(I,\varphi)\in\Lambda_1$, le foncteur composé
\diagram{\begin{tikzcd}[ampersand replacement=\&]
I_\bullet\arrow[r,"\varphi"] \& C\arrow[r,"f"] \& \mathcal N
\end{tikzcd}}
est $I$-\unsure{exact|ex.ct|coex.|coexact} \inserted{\unsure{à g.|à G}}, i.e. fait de $(f\varphi)(e_I)$ une $\underset{I}{\varinjlim}$ de $f\varphi|I$.
\scanpage{459}\mspage{371}
% archive sheet number: 458
b) \unsure{pour tout} $(I,\psi)\in\Lambda_2$, le foncteur composé
\diagram{\begin{tikzcd}[ampersand replacement=\&]
{}_\bullet I\arrow[r,"\psi"] \& C\arrow[r,"f"] \& \unsure{\mathcal N|N|\mathscr N}
\end{tikzcd}}
est $I$-exact à droite, i.e. fait de $(f\psi)(e_I)$ une limite projective de $(f\psi)|I$.

\deleted{\unsure{Sous}}\unsure{Avec} ces \deleted{conditions,} notations : si $\mathcal N$ est une catégorie accessible, il en est de même de $\underline{\mathrm{Hom}}_{\Lambda_1}^{\Lambda_2}(C,\mathcal N)$.

\uline{Commentaire} On utilisera ce théorème surtout quand on a $\Lambda_2=\varnothing$ ou $\Lambda_1=\varnothing$, auquel cas on note $\underline{\mathrm{Hom}}_{\Lambda_1}(C,\mathcal N)$, resp. $\underline{\mathrm{Hom}}^{\Lambda_2}(C,\mathcal N)$, la catégorie $\underline{\mathrm{Hom}}_{\Lambda_1}^{\Lambda_2}(C,\mathcal N)$. \deleted{De plus,} On appliquera le th. surtout quand $C$ elle-même est stable par $\unsure{\varinjlim|\underline{\lim}}$ de type $\Phi$ et \deleted{\unsure{par la}} $\varprojlim$ de type $\Psi$, et quand \deleted{les} $\Lambda_1$ (resp $\Lambda_2$) est déjà formé de foncteurs $\varphi$ (resp $\psi$) qui sont $I$-exacts à g. (resp. à droite). Donc $\underline{\mathrm{Hom}}_{\Lambda_1}^{\Lambda_2}(C,\mathcal N)$ est alors formée des $f:C\unsure{-|\longrightarrow}M$ qui commutent à certaines $\varinjlim$ de type $\Phi$, et à certaines $\varprojlim$ de type $\Psi$. En particulier, \unsure{quand}
\scanpage{460}\mspage{372}
% archive sheet number: 459
\uline{Corollaire} 4.14.2. \deleted{Avec les notations de} \inserted{Soient $C$ une catégorie} équivalente à une petite catégorie, $\unsure{\mathcal N|N|\mathscr N}$ une \inserted{$\unsure{\mathcal U|U|\mathscr U}$-}catégorie \deleted{accessible}, $\Phi$ et $\Psi$ deux \inserted{petits} ens. de catégories \inserted{$I$}, équivalentes à des petites catégories. On suppose \deleted{\unsure{stables}} $C$ \inserted{et $\mathcal N$} stables par $\varinjlim$ de type $\Phi$, par $\varprojlim$ de type $\Psi$. Soit $\underline{\mathrm{Hom}}_{\Phi}^{\Psi}(C,\mathcal N)$ la sous-catégorie str. pleine de $\underline{\mathrm{Hom}}(C,\mathcal N)$ formée des $f:C\longrightarrow M$ qui commutent aux $\varinjlim$ de type $\Phi$ et aux $\varprojlim$ de type $\Psi$. Alors, si $\mathcal N$ est accessible, il en est de même de $\underline{\mathrm{Hom}}_{\Phi}^{\Psi}(C,\mathcal N)$.

Voici des cas particuliers \deleted{\unsure{plus|bien}}\unsure{remarquables} de ce corollaire :

\uline{Corollaire} 4.14.3. Soient $M,N$ deux catégories accessibles, \inserted{et} $\pi_0$ un cardinal infini, tel que $M$, \deleted{\unsure{et}} $N$ satisfassent à $\unsure{\mathrm L_{\pi_0}|\mathrm L\pi_0}$. \unsure{Soient|Soit} $\underline{\mathrm{Hom}}_{\unsure{!|f|\xi}\pi_0}(M,N)$ la sous-catégorie \unsure{strictmt|strictement} pleine de $\underline{\mathrm{Hom}}(M,N)$ formée des foncteurs $\pi_0$-accessibles. Alors $\underline{\mathrm{Hom}}_{!\pi_0}(M,N)$ est accessible.
\scanpage{461}\mspage{373}
% archive sheet number: 460 (upper left)
(Le cas où $\pi_0$ est \unsure{utile} pour $M$ a déjà été traité.)

\uline{Dém} \deleted{\unsure{Supp-}} \deleted{\unsure{dans}} \inserted{Dans le cas général, soit} $\pi\unsure{>|\geq}\pi_0$ un cardinal \deleted{tel} utile pour $M$, de sorte qu'on a
\[
M\xleftarrow{\approx}\mathrm{Ind}_{\pi}(C)\qquad\text{avec}\quad C=M_{\pi}.
\]
\unsure{OPS} $M=\mathrm{Ind}_{\pi}(C)$. \deleted{\unsure{Alo-}} Le fait que $M$ satisfasse à $\unsure{\mathrm L_{\pi_0}|\mathrm L\pi_0}$ implique que $C$ est stable par $\unsure{\underline{\lim}|\varinjlim}$ gdes devant $\pi_0$, et de cardinal $\leq\pi$. Soit $\Phi$ l'ens des catégories de cette nature. \deleted{\unsure{$\xi$}} $\Phi$ n'est pas petit, mais il existe $\Phi_0\subset\Phi$ \deleted{\unsure{stable}} petit, tel que tout $I\in\Phi$ soit isom. à un $I_0\in\Phi_0$, ce qui \unsure{suffira|suffit} pour la validité de 4.14.1 et 4.14.2 (où on fera ici $\Psi=\varnothing$). On \deleted{peut} peut considérer $\underline{\mathrm{Hom}}_{!\pi_0}(M,N)$ comme une \deleted{\unsure{Hom}} sous-catégorie str. pleine de $\underline{\mathrm{Hom}}_{!\pi}(M,N)\xrightarrow{\approx}\underline{\mathrm{Hom}}(C,N)$ (cf. \unsure{4.13.3|4.17.3} p. 275). Pour un foncteur $f:M\longrightarrow N$ qui est $\pi$-accessible, dire qu'il est \inserted{$=$} $\pi_0$-accessible revient à dire que sa restriction \unsure{à} $C$, \unsure{i.e.|ou} son image dans $\underline{\mathrm{Hom}}(C,N)$, \deleted{est} commute aux $\varinjlim$ de type $\Phi$. \unsure{(cf.|Cf.} 4.13.20
% slanted
\margin{La référence n'est pas tout à fait \unsure{adéquate} : \unsure{il faudrait|il faut} \unsure{préciser} \unsure{ici} \unsure{que|les} \unsure{pour les} $\varinjlim$ de type $\Phi$ \unsure{il y a|on a} toujours les \unsure{conditions} \unsure{de petitesse|particulières} \unsure{requises|requis} ; \unsure{en termes|en toutes} \unsure{individuels|indiqués}.} page 341, où on prend $\alpha=\alpha_{\pi}$, noter que $I_{\alpha_\pi}$ est gde devant $\pi_0$ puisque $\pi>\pi_0$.)
\scanpage{462}\mspage{374}
% archive sheet number: 461
Donc on trouve
\[
\underline{\mathrm{Hom}}_{!\pi_0}(M,N)\xrightarrow{\approx}\underline{\mathrm{Hom}}_{\Phi}(C,N),
\]
et le corollaire 4.14.2. \unsure{donne} la conclusion voulue.

\uline{Corollaire} 4.14.4. Soient $M,N$ deux catégories accessibles, $\Phi$ et $\Psi$ deux ens. de catégories comme dans \unsure{4.14.1|4.14.2}, \deleted{\unsure{$N$} \illegible{1} \unsure{catégorie}} \inserted{on suppose que} \inserted{$M$ et} \deleted{\unsure{$C$}} $N$ \deleted{est} \inserted{satisfont $\unsure{\mathrm L_{\pi_0}|\mathrm L\pi_0}$, et \deleted{\unsure{sont}}} stables par $\varinjlim$ de type $\Phi$ et par $\varprojlim$ de type $\Psi$, Soit $\underline{\mathrm{Hom}}_{\Phi,\pi_0}^{\Psi}(M,N)$ la sous-catégorie str. pl. de $\underline{\mathrm{Hom}}_{\deleted{\unsure{\pi}}}(M,N)$ \unsure{formée|formé} des foncteurs $\pi_0$-accessibles, \deleted{et tels que} qui commutent aux $\varinjlim$ de type $\Phi$, et aux $\varprojlim$ de type $\Psi$. Alors $\underline{\mathrm{Hom}}_{\Phi,\pi_0}^{\Psi}(M,N)$ est accessible.

\uline{Dém} : Cette fois, on choisit $\pi>\pi_0$ tel qu'on ait $\operatorname{card}\unsure{I_{\mathrm{red}}^{.}|I\mathrm{red}^{.}|I_{\mathrm{red}}|I\mathrm{red}}\leq\unsure{\pi}$ pour tt $I$ dans $\Phi$ ou dans $\Psi$. La catégorie envisagée s'interprète alors (posant encore $C=M_{\pi}$) comme \deleted{\unsure{$K_{\alpha}$|la}}
\scanpage{463}\mspage{375}
% archive sheet number: 462
\deleted{\unsure{$X$}}
\[
\underline{\mathrm{Hom}}_{\Phi,\pi_0}^{\Psi}(M,N)\xrightarrow{\approx}
\underline{\mathrm{Hom}}_{\Phi'}^{\Psi}(C,N),
\]
où $\Phi'$ est \unsure{formée|formé} des catégories qui sont \deleted{soit dans} $\Phi$, ou qui sont gdes devant $\pi_0$ et de cardinal $\leq\pi$. Il suffit \unsure{donc} \deleted{\unsure{de}}d'appliquer encore \unsure{4.14.2|4.14.?}.

\uline{Corollaire} 4.14.5. \deleted{Soit $\Phi$ l'un des} \inserted{Considérons les} trois ensembles suivants : a) $\Phi_1$ est formée de toutes les catégories \unsure{*} équivalentes à des petites catégories \quad b) \deleted{On suppose \unsure{toujours}} \inserted{$\Phi_2$ comme a), mais en} se restreignant aux catégories filtrantes.

% slanted
\margin{On pourrait \deleted{\unsure{en fait|en tout}} faire ici un énoncé qui englobe 4.14.4, en rajoutant un d) $\Phi_4$ formé des catégories gdes devant $\pi_0$ et prendre $\Phi$ \deleted{de la forme \unsure{$\Phi$} \unsure{de}} \inserted{\unsure{de la forme|l'un des}} $\Phi_1,\deleted{\unsure{\Phi_2}}\Phi_2,\Phi_4,\Phi_3,$ \deleted{\unsure{ou former}} $\Phi_3\cup\Phi_2,\Phi_3\cup\Phi_4$ \deleted{\unsure{avec ou} $\Phi_2$ \unsure{ou} $\Phi_4$ ou $\Phi_3,\Phi_2$} ou bien la réunion d'un de ceux-ci et d'un petit ensemble $\Phi_0$ \unsure{de|des} cat. \unsure{ess. petites|ens. petites|$\mathcal U$-petites}.}

c) $\Phi_3$ \deleted{est} formé des catégories discrètes, \unsure{de|ou} celles qui sont gdes devant $\pi_0$ ($\pi_0$ un cardinal infini donné). Soit $\Phi$ \deleted{un \unsure{ensemble} de catégories, \unsure{soumises} \unsure{ici}} \inserted{l'un des ens. de catégories} $\Phi_1,\Phi_2,\Phi_3,\Phi_2\cup\Phi_3$, \unsure{;} soit $\Psi$ un \inserted{petit} ensemble de catégories. Soient $M,N$ \deleted{deux} \unsure{catégorie|catégories} \unsure{accessibles}, stables par $\underset{\deleted{\unsure{\gets}}\to}{\lim}$ de type $\Phi$, et par $\varprojlim$ de type $\Psi$. Alors \deleted{pour tte catégorie accessible $N$,} $\underline{\mathrm{Hom}}_\Phi^\Psi(M,N)$ est accessible.
\scanpage{464}\mspage{376}
% archive sheet number: 463
\uline{Démonstration} Soit $\pi'_0$ \deleted{\unsure{un cardinal}} \inserted{$\unsure{\geq|>}\pi_0$ un} cardinal utile pour $M$, $\pi$ un cardinal ayant les propriétés suivantes
\[
\begin{aligned}
\text{a)}\quad &\pi\geq\operatorname{Sup}(\pi'_0,\operatorname{card\,red}M_{\pi'_0}),\quad\pi>\pi_0\\
\text{b)}\quad &\pi^{\pi'_0}=\pi\\
\text{c)}\quad &\underline I\in\deleted{\unsure{\Phi\cup\Psi}}\inserted{\Phi_0\cup\Psi}\Longrightarrow
\operatorname{card\,red}\underline I\leq\pi .
\end{aligned}
\]
\deleted{\unsure{$\mathrm{avec}\ \Phi$} \illegible{2} \unsure{$I$}} \deleted{On voit} Soit $C=M_\pi$, d'où $M\underset{\unsure{-|\leftarrow}}{\approx}\mathrm{Ind}_\pi(C)$. Dans tous les cas envisagés, $\Phi$ contient l'ens des cat. \inserted{\unsure{ess.|ens.|équ.} petites} qui sont gdes devant $\pi_0$, \unsure{a|à} fortiori celles qui sont gdes devant $\pi$, donc la catégorie $\underline{\mathrm{Hom}}_\Phi^\Psi$ \deleted{\unsure{est}} \inserted{s'interprète comme} une sous-catégorie \unsure{strictmt|strictement} pleine de
\[
\underline{\mathrm{Hom}}(C,N).
\]
Si d'autre part $\Phi_\pi$ désigne l'ens des $I\in\Phi$ telles que $\operatorname{card} I\leq\unsure{\pi|1}$, on trouve
\[
\underline{\mathrm{Hom}}_\Phi^\Psi(M,N)\underset{\unsure{-|\to}}{\approx}
\underline{\mathrm{Hom}}_{\Phi_\pi}^\Psi(C,N),
\]
et on applique 4.14.2.
\scanpage{465}\mspage{377}
% archive sheet number: 464
En particulier, on trouve

\uline{Corollaire} 4.14.6. Soient $M,N$ deux catégories accessibles, supposons $M$ stable par petites $\varinjlim$ (i.e $M$ est une cat. \unsure{paratopos}). Alors la catégorie $\underline{\mathrm{Hom}}_{\deleted{\unsure{\Phi}}!}(M,N)$ des foncteurs $f:M\longrightarrow N$ qui commutent aux \inserted{petites} $\varinjlim$ est accessible. Si $N$ est elle aussi un paratopos, $\underline{\mathrm{Hom}}_!(M,N)$ est un paratopos.

\uline{Dém} La première assertion est le cas particulier de 4.14.5 relatif au cas où $\Phi=\Phi_0$, $\Psi=\varnothing$. La deuxième assertion en résulte, puisque $\underline{\mathrm{Hom}}_!(M,N)$ est \unsure{visiblemt|visiblement} stable par petites $\varinjlim$, \unsure{donc|car} c'est \unsure{un|une} paratopos étant accessible (4.8.9 page 207).

\uline{Corollaire} 4.14.7. Soient $\Phi,\Psi,$\inserted{$M$} comme dans 4.14.5, \deleted{et soit $M$ une catégorie accessible telle que $M^\circ$ (catégorie opposée à $M$) soit stable par $\varinjlim$ de type $\Phi$, par $\varprojlim$ de type $\Psi$, i.e. telle que} \inserted{Alors pour \deleted{\unsure{toute}} catégorie accessible $N$, la catégorie}
\scanpage{466}\mspage{378}
% archive sheet number: 465
\deleted{$M$ soit stable par $\varinjlim$ \uline{de type} $\Psi$, par $\varprojlim$ de type $\Phi$. Soit d'autre part $N$ une catégorie accessible, et considérons la \deleted{catégorie}}
\[
\underline{\mathrm{Hom}}_{\deleted{\unsure{\Phi}}\Psi}^{\deleted{\unsure{\Psi}}\Phi}(M^\circ,N)
\underset{\unsure{-|\to}}{\sim}
\left(\underline{\mathrm{Hom}}_{\Phi}^{\Psi}(M,N^\circ)\right)^\circ
\]
est accessible.

\uline{Dém} On procède comme pour 4.14.5, pour expliciter
\[
\underline{\mathrm{Hom}}_{\Phi}^{\Psi}(M,N^\circ)\xrightarrow{\approx}
\underline{\mathrm{Hom}}_{\Phi_\pi}^{\Psi}(C,N^\circ)
\]
d'où
\[
\underline{\mathrm{Hom}}_{\Psi}^{\Phi}(M^\circ,N)
\underset{\unsure{\leftrightarrow|\leftarrow|\to}}{\approx}
\underline{\mathrm{Hom}}_{\Psi}^{\Phi_\pi}(C^\circ,N),
\]
et on applique 4.14.2 \unsure{à} $C^\circ,\Psi,\Phi_\pi$.

\uline{Corollaire} 4.14.8. Soi\deleted{\unsure{t}}\unsure{ent} $M$ et $N$ \unsure{deux|des} catégories \deleted{accessibles, supposons \deleted{\unsure{par tt}} $M$ stable par petites $\varinjlim$. Alors la catégorie} \inserted{paratopos. Alors}
\[
M\boxtimes N\underset{\unsure{-|\to}}{\approx}\underline{\mathrm{Hom}}^!(M^\circ,N)
\qquad\text{(cf 4.2 et 4.3)}
\]
est un paratopos.

En effet, en vertu de 4.14.7 (\unsure{avec|cas} $\Phi=\Phi_{\unsure{M|1|0}}$, $\Psi=\varnothing$) c'est une catégorie accessible.
\scanpage{467}
\mspage{379}
% archive sheet number: 466

D'autre part il est immédiat qu'elle est stable par \deleted{\unsure{pe}}petites $\varprojlim$, $N$ l'étant. Donc on conclut \deleted{\unsure{par}} \unsure{encore|en} par 4.8.9 (page 207) que c'est une catégorie paratopos.

\uline{Corollaire} 4.14.9. Soient $M,N$ deux catégories paratopos, et $\underline{\mathrm{Hom}}^{!}(M,N)$ la catégorie des foncteurs $M\longrightarrow N$ qui commutent aux \inserted{petites} $\varprojlim$. On a alors une équivalence de catégories
\[
\begin{aligned}
\underline{\mathrm{Hom}}^{!}(M,N)&\underset{\unsure{-|\leftarrow}}{\approx}\underline{\mathrm{Hom}}_{!}(N,M)^{\circ}\\
\deleted{\text{\unsure{ce qui est|qui associe} \unsure{à tout}}}\qquad f&\longmapsto\text{adjoint à gauche de }f,
\end{aligned}
\]
\margin{\figure{\begin{tikzpicture}[x=.0058cm,y=-.0058cm,line cap=round,line join=round]
\draw[line width=.65pt] (2,30) .. controls (45,1) and (110,-3) .. (149,1) .. controls (202,3) and (217,40) .. (195,83) .. controls (176,123) and (123,156) .. (73,196) .. controls (32,226) and (-2,266) .. (16,286) .. controls (28,302) and (64,304) .. (91,298) .. controls (141,291) and (184,285) .. (229,272);
\draw[line width=.8pt] (2,107) .. controls (4,124) and (4,149) .. (14,167);
\draw[line width=.9pt] (1,176) -- (8,184) -- (11,181);
\draw[line width=.8pt] (244,87) .. controls (247,107) and (246,132) .. (252,149);
\draw[line width=1pt] (254,166) -- (258,176);
\end{tikzpicture}}}%
et par suite la catégorie \underbar{\uuline{opposée}} à $\underline{\mathrm{Hom}}^{!}(M,N)$ est accessible.

Ceci résulte du

\uline{Lemme} 4.14.10 (pour mémoire) \unsure{Soient|Soit} $M$ \deleted{\unsure{et}} \deleted{\unsure{une}} cat.-paratopos, $f:M\longrightarrow N$ un foncteur. Pour que $f$ admette un adjoint à g. (resp à droite) il f. et s. que $f$ commute aux petites $\varprojlim$ \unsure{(|((}resp aux petites $\varinjlim$).
